% Append inside the root paper's thebibliography environment, adjusting keys
% if the root already provides these references.
\bibitem{KahnHeight2024}
B. Kahn, with an appendix by Q. Liu,
\emph{Refined height pairing}, Algebra \& Number Theory \textbf{18}
(2024), no. 6, 1039--1079,
doi:10.2140/ant.2024.18.1039.
Locators: Definition 2.2; Proposition 2.8; Definition 2.10 and (2-9);
Propositions 3.6, 3.8, 3.9 and 4.12; Theorem 4.14(b);
Theorem 5.6(a); for the classical sign,
Theorem 5.10(b).
\url{https://msp.org/ant/2024/18-6/ant-v18-n6-p01-p.pdf}.

\bibitem{ChenHeight2025}
Z. Chen, \emph{Degeneration of the archimedean height pairing of
algebraically trivial cycles}, arXiv:2512.22788v1,
submitted 28 December 2025; downloaded PDF dated 30 December 2025.
Locators: Lemma 3.16; Definition 3.17; Lemma 3.19; Example 3.20;
Remark 3.21 (perfect-field descent);
Theorems 5.10, 5.13, 7.5, 7.9, and 7.19(a),(b).
The Neron--Tate comparison in 7.9 is distinct from the further
Griffiths-conditional Lear comparison in 7.19(c).
\url{https://arxiv.org/pdf/2512.22788v1}.

\bibitem{KunnemannPicard1996}
K. Kunnemann, \emph{Higher Picard varieties and the height pairing},
American Journal of Mathematics \textbf{118} (1996), no. 4,
781--797, doi:10.1353/ajm.1996.0033.  Locators: Lemma 8.1 and
Theorem 8.2.  Cited as the number-field antecedent of Chen's
geometric Theorem 7.9; no extension of the number-field theorem
to function fields is silently assumed.

\bibitem{SaitoPicard1979}
H. Saito, \emph{Abelian varieties attached to cycles of intermediate
dimension}, Nagoya Mathematical Journal \textbf{75} (1979),
95--119, doi:10.1017/S0027763000024855.
Locators: Theorem 3.6; (4.1); Proposition 4.4; Corollary 4.5.
The ground field is explicitly algebraically closed of arbitrary
characteristic (p.~96).  The Poincare-cycle and incidence-duality
formulas are restated in Chen, Lemmas 3.16 and 3.19.

\bibitem{PepinNeron2011}
C. P\'epin, \emph{Neron's pairing and relative algebraic equivalence},
arXiv:1103.0570v1 (2011), Definition 2.1 and Theorem 2.3,
proof pp.~6--8; see also the curve intersection formula on p.~2.
The base is a complete discrete valuation ring with algebraically
closed residue field, without a characteristic restriction.
\url{https://arxiv.org/pdf/1103.0570}.

\bibitem{SGA7IHeight}
A. Grothendieck, M. Raynaud and D.~S. Rim,
\emph{Groupes de monodromie en geometrie algebrique, I},
SGA 7 I, Lecture Notes in Mathematics \textbf{288}, Springer, 1972.
Expose IX: Corollary 2.2.9, pp.~335--336 (good reduction);
Proposition 3.2 and 3.2.1, pp.~347--348 (global Neron model on a
regular curve and quasi-projectivity).  In the retrieved scan these
are PDF pages 340--341 and 352--353, respectively; the expose's
internal page numbers are 17--18 and 29--30.
Original scan retrieved from
\url{https://library.slmath.org/nonmsri/sga/sga/pdf/sga7-1.pdf}.

% Existing shared bibliography keys to reuse if present:
% Fulton: W. Fulton, Intersection Theory, 2nd ed., 1998,
% Proposition 1.8; Corollary 6.3; Proposition 6.7(e);
% Example 6.3.4(a) / Proposition 14.1(c).
% Deligne Weil I: Publ. Math. IHES 43 (1974), 273--307,
% (4.1), (4.3.3), (5.4), (5.8)(a),(b).
% SGA 7 II, LNM 340 (1973), Expose XIV 3.2; Expose XVIII 6.6.
